% !TEX root = ../main.tex
%% Corrigés — Chapitre 18 — Espaces préhilbertiens réels
%% Résultats vérifiés par verifications/chapitre_18.py.

\corrigesChapitre{18}{Espaces préhilbertiens réels}
\corrigesSouples

\begin{blocLivre}{cor-18-1}
\begin{corrige}{1}{Produit scalaire ou pas ?}
\demarche{Bilinéaire et symétrique se voient ; tout se joue sur « défini
positif ».}
\begin{questions}
  \item $\varphi$ : \res{\text{non}} — $\varphi(u,u)=x^{2}-y^{2}$ peut
    être négatif ($u=(1,2)$ donne $-3$). $\psi$ : \res{\text{oui}} —
    bilinéaire symétrique, et $\psi(u,u)=2x^{2}+3y^{2}>0$ pour $u\neq0$.
  \item $\int_0^1fg$ : \res{\text{oui}} — le caractère défini vient de
    l'exercice « intégrale nulle » du chapitre 15. $f(0)g(0)$ :
    \res{\text{non}} — $f(x)=x$ est non nulle avec
    $\langle f,f\rangle=0$ : positive mais pas \emph{définie}.
\end{questions}
\end{corrige}
\end{blocLivre}

\begin{blocLivre}{cor-18-2}
\begin{corrige}{2}{Calculs dans $\R^{3}$}
\demarche{Produit scalaire canonique : sommes de produits, racines pour
les normes.}
\begin{questions}
  \item $\norme{u}=\sqrt{1+4+4}=\res{3}$, $\norme{v}=3$ ;
    $\langle u,v\rangle=2-2+4=\res{4}$, d'où
    $\cos\theta=\dfrac{4}{3\times3}=\res{\dfrac49}$.
  \item $u-v=(-1,3,0)$ : $\norme{u-v}=\res{\sqrt{10}}$.
\end{questions}
\end{corrige}
\end{blocLivre}

\begin{blocLivre}{cor-18-3}
\begin{corrige}{3}{Cauchy-Schwarz en pratique}
\demarche{Choisir le bon second vecteur : ici (1, 1, ..., 1).}
\begin{questions}
  \item C'est Cauchy-Schwarz pour $a=(a_1,\dots,a_n)$ et
    $y=(1,\dots,1)$ :
    $\langle a,y\rangle^{2}\leqslant\norme{a}^{2}\norme{y}^{2}
    =n\sum a_i^{2}$. Égalité si et seulement si $a$ est colinéaire à
    $y$ : \res{\text{tous les }a_i\text{ égaux}}.
  \item $(1+2+3+4)^{2}=100\leqslant4\times30=120$ ; et pour $(2,2,2)$ :
    $36=3\times12$ — égalité, comme prévu.
\end{questions}
\end{corrige}
\end{blocLivre}

\begin{blocLivre}{cor-18-4}
\begin{corrige}{4}{Polarisation et parallélogramme}
\demarche{Développer les carrés de normes : tout est identité remarquable.}
\begin{questions}
  \item $\langle u,v\rangle=2+3-2=3$ ; $u+v=(3,4,1)$ :
    $\norme{u+v}^{2}=26=6+14+2\times3$ : conforme.
  \item $u-v=(1,-2,-3)$, $\norme{u-v}^{2}=14$ :
    $\dfrac{26-14}{4}=\res{3}=\langle u,v\rangle$ (polarisation), et
    $26+14=40=2\times6+2\times14$ (parallélogramme).
\end{questions}
\end{corrige}
\end{blocLivre}

\begin{blocLivre}{cor-18-5}
\begin{corrige}{5}{Orthogonalité et Pythagore}
\demarche{Pour la liberté : prendre le produit scalaire de la relation avec
chaque vecteur.}
\begin{questions}
  \item $\langle u,v\rangle=-2+2=0$ ; et
    $\norme{u+v}^{2}=\norme{(-1,3)}^{2}=10=5+5$ : Pythagore.
  \item Si $\sum\lambda_ix_i=0$, le produit scalaire avec $x_j$ donne
    $\lambda_j\norme{x_j}^{2}=0$ (les autres termes s'annulent par
    orthogonalité), donc $\lambda_j=0$ : la famille est
    \res{\text{libre}}.
\end{questions}
\end{corrige}
\end{blocLivre}

\begin{blocLivre}{cor-18-6}
\begin{corrige}{6}{Gram-Schmidt dans $\R^{3}$}
\demarche{Retrancher les ombres, normer à la fin.}
$e_1=\dfrac{1}{\sqrt2}(1,1,0)$. Puis
$w_2=(1,0,1)-\dfrac12(1,1,0)=\bigl(\tfrac12,-\tfrac12,1\bigr)$, de norme
$\sqrt{3/2}$ : $e_2=\dfrac{1}{\sqrt6}(1,-1,2)$. Enfin
\[
w_3=(0,1,1)-\frac12(1,1,0)-\frac16(1,-1,2)
=\Bigl(-\frac23,\frac23,\frac23\Bigr),
\qquad
e_3=\res{\frac{1}{\sqrt3}(-1,1,1)} .
\]
On vérifie : les trois produits scalaires croisés sont nuls, les trois
normes valent $1$.
\end{corrige}
\end{blocLivre}

\begin{blocLivre}{cor-18-7}
\begin{corrige}{7}{Coordonnées en base orthonormée}
\demarche{En BON, les coordonnées sont des produits scalaires.}
\begin{questions}
  \item $\langle x,e_1\rangle=\dfrac{3+8}{5}=\res{\dfrac{11}{5}}$ et
    $\langle x,e_2\rangle=\dfrac{-4+6}{5}=\res{\dfrac25}$, et l'on
    vérifie $\frac{11}{5}e_1+\frac25e_2=(1,2)$.
  \item $\Bigl(\dfrac{11}{5}\Bigr)^{2}+\Bigl(\dfrac25\Bigr)^{2}
    =\dfrac{125}{25}=5=\norme{x}^{2}$ : la norme se lit dans n'importe
    quelle BON.
\end{questions}
\end{corrige}
\end{blocLivre}

\begin{blocLivre}{cor-18-8}
\begin{corrige}{8}{Projection sur une droite}
\demarche{Projeter sur une droite : un seul produit scalaire.}
\begin{questions}
  \item $p(x)=\dfrac{\langle x,u\rangle}{\norme{u}^{2}}\,u
    =\dfrac39(1,2,2)=\res{\Bigl(\dfrac13,\dfrac23,\dfrac23\Bigr)}$.
  \item $d(x,D)^{2}=\norme{x}^{2}-\norme{p(x)}^{2}=9-1=8$ :
    $d(x,D)=\res{2\sqrt2}$, et $9=1+8$ : Pythagore au rendez-vous.
\end{questions}
\end{corrige}
\end{blocLivre}

\begin{blocLivre}{cor-18-9}
\begin{corrige}{9}{Le produit scalaire des matrices}
\demarche{$\operatorname{tr}\bigl(A^{\mathsf T}B\bigr)$ additionne les
produits coefficient par coefficient : la géométrie de $\R^{4}$ en
déguisement.}
\begin{questions}
  \item $\bigl(A^{\mathsf T}B\bigr)_{jj}=\sum_ia_{ij}b_{ij}$ : en
    sommant, $\langle A,B\rangle=\sum_{i,j}a_{ij}b_{ij}$ ; ici
    $0+2+3-2=\res{3}$.
  \item $\langle E_{ij},E_{kl}\rangle=1$ si $(i,j)=(k,l)$, $0$ sinon :
    les quatre matrices élémentaires forment une \res{\text{BON}} de
    $\mathcal{M}_2(\R)$.
  \item Le trinôme du chapitre 13 était la démonstration de
    Cauchy-Schwarz pour ce produit scalaire — énoncée alors « avant
    l'heure », désormais un cas particulier du cours.
\end{questions}
\end{corrige}
\end{blocLivre}

\begin{blocLivre}{cor-18-10}
\begin{corrige}{10}{Le produit scalaire intégral}
\demarche{Les fonctions deviennent des vecteurs ; les angles ont un sens.}
\begin{questions}
  \item Bilinéarité et symétrie sont claires ; si
    $\int_0^1f^{2}=0$ avec $f$ continue, alors $f=0$ (chapitre 15,
    exercice 10) : défini positif.
  \item $\norme{x}^{2}=\int_0^1t^{2}=\dfrac13$,
    $\langle1,x\rangle=\dfrac12$, et
    $\cos\theta=\dfrac{1/2}{1\times1/\sqrt3}=\res{\dfrac{\sqrt3}{2}}$ :
    les fonctions $1$ et $t\mapsto t$ font un angle de $\dfrac\pi6$.
\end{questions}
\end{corrige}
\end{blocLivre}

\begin{blocLivre}{cor-18-11}
\begin{corrige}{11}{Gram-Schmidt sur les polynômes}
\demarche{Le même algorithme, avec des intégrales pour produits scalaires.}
\begin{questions}
  \item $P_1=X-\dfrac{\langle X,1\rangle}{\langle1,1\rangle}
    =X-\dfrac12$ ; puis
    $P_2=X^{2}-\dfrac{\langle X^{2},1\rangle}{1}
    -\dfrac{\langle X^{2},P_1\rangle}{\norme{P_1}^{2}}P_1
    =\res{X^{2}-X+\dfrac16}$.
  \item $\norme{P_1}^{2}=\dfrac{1}{12}$ et
    $\norme{P_2}^{2}=\dfrac{1}{180}$ ; et
    $6P_2=6X^{2}-6X+1$ : le polynôme orthogonal à $1$ et $X$ de
    l'exercice 22 du chapitre 15 — c'était un Gram-Schmidt qui
    s'ignorait.
\end{questions}
\end{corrige}
\end{blocLivre}

\begin{blocLivre}{cor-18-12}
\begin{corrige}{12}{L'orthogonal d'un sous-espace}
\demarche{Une équation = un vecteur normal : l'orthogonal se lit
directement.}
\begin{questions}
  \item $x\perp(1,1,1)$ signifie $x+y+z=0$ : $F^{\perp}$ est le plan
    d'équation $x+y+z=0$, de dimension $\res{2}$, de base
    $\bigl((1,-1,0),(1,0,-1)\bigr)$.
  \item $F\cap F^{\perp}=\ens{0}$ (un vecteur orthogonal à lui-même est
    nul) et $1+2=3$ : somme directe. Enfin
    $F\subseteq(F^{\perp})^{\perp}$ et les dimensions ($1=3-2$)
    forcent l'égalité : $\res{(F^{\perp})^{\perp}=F}$.
\end{questions}
\end{corrige}
\end{blocLivre}

\begin{blocLivre}{cor-18-13}
\begin{corrige}{13}{Projection sur un plan}
\demarche{Retirer la composante normale : un produit scalaire suffit.}
\begin{questions}
  \item $n=(1,1,1)$ et
    $p_H(a)=a-\dfrac{\langle a,n\rangle}{\norme{n}^{2}}\,n
    =(1,2,3)-2(1,1,1)=\res{(-1,0,1)}$.
  \item $d(a,H)=\norme{(2,2,2)}=\res{2\sqrt3}$. La décomposition
    $(1,2,3)=(-1,0,1)+(2,2,2)$ du chapitre 12 était donc
    \emph{orthogonale} : on sait désormais que $(-1,0,1)$ n'est pas un
    supplémentaire parmi d'autres, mais le point du plan \emph{le plus
    proche} de $a$.
\end{questions}
\end{corrige}
\end{blocLivre}

\begin{blocLivre}{cor-18-14}
\begin{corrige}{14}{Distance à un hyperplan}
\demarche{Décomposer x le long du vecteur normal.}
\begin{questions}
  \item Écrire $x=p+\lambda u$ avec $p\in H$ : le produit scalaire par
    $u$ donne $\lambda=\dfrac{\langle x,u\rangle}{\norme{u}^{2}}$, et
    $d(x,H)=\norme{\lambda u}
    =\res{\dfrac{|\langle x,u\rangle|}{\norme{u}}}$.
  \item $u=(1,2,-2)$, $\norme{u}=3$ :
    $d=\dfrac{|3+2-2|}{3}=\res{1}$.
\end{questions}
\end{corrige}
\end{blocLivre}

\begin{blocLivre}{cor-18-15}
\begin{corrige}{15}{Le projeté en base orthonormée}
\demarche{Deux produits scalaires, deux coordonnées : le projeté tombe.}
\begin{questions}
  \item Avec la BON $e_1=\frac{1}{\sqrt2}(1,1,0)$,
    $e_2=\frac{1}{\sqrt6}(1,-1,2)$ de $F$ (exercice 6) :
    $\langle x,e_1\rangle=\frac{3}{\sqrt2}$,
    $\langle x,e_2\rangle=\frac{3}{\sqrt6}$, d'où
    \[
    p_F(x)=\frac32(1,1,0)+\frac12(1,-1,2)=\res{(2,1,1)} .
    \]
  \item $x-p_F(x)=(1,-1,-1)$ : $d(x,F)=\res{\sqrt3}$, et l'on vérifie
    $\langle(1,-1,-1),e_1\rangle=\langle(1,-1,-1),e_2\rangle=0$ — le
    reste est bien orthogonal au plan, et $9=6+3$.
\end{questions}
\end{corrige}
\end{blocLivre}

\begin{blocLivre}{cor-18-16}
\begin{corrige}{16}{La constante la plus proche}
\demarche{Minimiser une distance quadratique, c'est projeter.}
\begin{questions}
  \item Le problème demande le point de
    $F=\operatorname{Vect}(1)$ le plus proche de $f$ : c'est
    $p_F(f)=\dfrac{\langle f,1\rangle}{\langle1,1\rangle}
    =\displaystyle\int_0^1f$ — la \res{\text{valeur moyenne}}
    (chapitre 15), revisitée en géomètre.
  \item Pour $f(t)=t^{2}$ : $c=\dfrac13$ et le minimum vaut
    $\displaystyle\int_0^1\Bigl(t^{2}-\frac13\Bigr)^{2}\mathrm{d}t
    =\res{\dfrac{4}{45}}$.
\end{questions}
\end{corrige}
\end{blocLivre}

\begin{blocLivre}{cor-18-17}
\begin{corrige}{17}{Moindres carrés polynomiaux}
\demarche{Le reste doit être orthogonal à 1 et à t : deux équations
linéaires.}
L'orthogonalité de $t^{2}-a-bt$ à $1$ puis à $t$ s'écrit
$\dfrac13=a+\dfrac b2$ et $\dfrac14=\dfrac a2+\dfrac b3$, d'où
$\res{a=-\dfrac16,\quad b=1}$ : le projeté de $X^{2}$ sur $\R_1[X]$ est
$X-\dfrac16$, et le minimum vaut
$\norme{X^{2}-X+\tfrac16}^{2}=\norme{P_2}^{2}=\res{\dfrac{1}{180}}$
(exercice 11) : la meilleure droite approchant la parabole sur
$\intFF{0}{1}$.
\end{corrige}
\end{blocLivre}

\begin{blocLivre}{cor-18-18}
\begin{corrige}{18}{Pythagore généralisé et inégalité de Bessel}
\demarche{Le projeté et le reste sont orthogonaux : Pythagore fait tout.}
\begin{questions}
  \item Développer $\norme{\sum x_i}^{2}$ : les doubles produits
    $\langle x_i,x_j\rangle$ sont nuls, il reste
    $\sum\norme{x_i}^{2}$.
  \item Posons $p=\sum_i\langle x,e_i\rangle e_i$ : alors $x-p$ est
    orthogonal à chaque $e_i$, donc à $p$, et Pythagore donne
    $\norme{x}^{2}=\norme{p}^{2}+\norme{x-p}^{2}
    \geqslant\norme{p}^{2}=\sum_i\langle x,e_i\rangle^{2}$ : c'est
    \res{\text{Bessel}}, avec égalité si et seulement si $x=p$,
    c'est-à-dire $x\in\operatorname{Vect}(e_1,\dots,e_p)$.
\end{questions}
\end{corrige}
\end{blocLivre}

\begin{blocLivre}{cor-18-19}
\begin{corrige}{19}{La meilleure approximation cubique}
\demarche{Trois conditions d'orthogonalité, un système linéaire — ou la
famille orthogonale toute prête.}
L'orthogonalité de $t^{3}-at^{2}-bt-c$ à $1$, $t$, $t^{2}$ donne un
système linéaire de solution
$a=\dfrac32$, $b=-\dfrac35$, $c=\dfrac{1}{20}$ : le projeté est
$\dfrac32X^{2}-\dfrac35X+\dfrac1{20}$, et
\[
\min=\int_0^1\Bigl(t^{3}-\frac32t^{2}+\frac35t-\frac1{20}\Bigr)^{2}
\mathrm{d}t=\res{\frac{1}{2800}} .
\]
Le polynôme d'erreur est le quatrième membre de la famille orthogonale —
exactement le $P_3$ de l'exercice 24.
\end{corrige}
\end{blocLivre}

\begin{blocLivre}{cor-18-20}
\begin{corrige}{20}{Les polynômes de Legendre}
\demarche{Sur un segment symétrique, la parité offre la moitié des
orthogonalités.}
\begin{questions}
  \item Les produits pair-impair sont nuls par imparité ; restent
    $\langle1,X^{2}-\tfrac13\rangle=\tfrac23-\tfrac23=0$ et
    $\langle X,X^{3}-\tfrac35X\rangle=\tfrac25-\tfrac35\cdot\tfrac23=0$.
  \item $\dfrac{\mathrm{d}^{2}}{\mathrm{d}x^{2}}(x^{2}-1)^{2}
    =12x^{2}-4=12\bigl(x^{2}-\tfrac13\bigr)$ et
    $\dfrac{\mathrm{d}^{3}}{\mathrm{d}x^{3}}(x^{2}-1)^{3}
    =120\bigl(x^{3}-\tfrac35x\bigr)$ : Rodrigues confirmé pour $n=2,3$.
  \item Une famille orthogonale avec $\deg P_n=n$ est déterminée à
    scalaires près : chaque $P_n$ est orthogonal à $\R_{n-1}[X]$ et de
    degré $n$, ce qui ne laisse qu'une direction — celle que
    Gram-Schmidt construit.
\end{questions}
\end{corrige}
\end{blocLivre}

\begin{blocLivre}{cor-18-21}
\begin{corrige}{21}{Le déterminant de Gram}
\demarche{Développer : Gram = Cauchy-Schwarz au carré, et c'est une aire.}
\begin{questions}
  \item $G(u,v)=\norme{u}^{2}\norme{v}^{2}-\langle u,v\rangle^{2}
    \geqslant0$ : c'est exactement \res{\text{Cauchy-Schwarz}}, et la
    nullité est son cas d'égalité : famille liée.
  \item $u=(3,1)$, $v=(1,2)$ :
    $G=10\times5-5^{2}=\res{25}=5^{2}$ — le carré de l'aire du
    parallélogramme du chapitre 14. En général,
    $G(u,v)=\norme{u}^{2}\norme{v}^{2}\sin^{2}\theta
    =(\text{base}\times\text{hauteur})^{2}$ : le déterminant de Gram
    mesure les aires \emph{sans coordonnées}, dans tout préhilbertien.
\end{questions}
\end{corrige}
\end{blocLivre}

\begin{blocLivre}{cor-18-22}
\begin{corrige}{22}{Familles obtusangles}
\demarche{Séparer les coefficients positifs et négatifs d'une relation :
les produits croisés travaillent tous dans le même sens.}
\begin{questions}
  \item Trois vecteurs du plan à $120$ degrés deux à deux :
    $\langle x_i,x_j\rangle=-\dfrac12<0$.
  \item Supposons $\sum_{i<p}\lambda_ix_i=0$ non triviale, et notons
    $y=\sum_{\lambda_i>0}\lambda_ix_i=\sum_{\lambda_j<0}(-\lambda_j)x_j$.
    Alors $\norme{y}^{2}$ est une somme de termes
    $\lambda_i(-\lambda_j)\langle x_i,x_j\rangle\leqslant0$ : $y=0$.
    L'un des deux camps est non vide, disons
    $\sum_{\lambda_i>0}\lambda_ix_i=0$ ; son produit scalaire avec
    $x_p$ est strictement négatif (termes $\lambda_i\langle
    x_i,x_p\rangle<0$) et nul à la fois : contradiction. La sous-famille
    est \res{\text{libre}}.
  \item $p$ vecteurs obtusangles donnent $p-1$ vecteurs libres :
    $p-1\leqslant n$, soit $\res{p\leqslant n+1}$ — borne atteinte par le
    « triangle » de a) et ses analogues en toute dimension.
\end{questions}
\end{corrige}
\end{blocLivre}

\begin{blocLivre}{cor-18-23}
\begin{corrige}{23}{L'inégalité de Hadamard}
\demarche{Gram-Schmidt rend les colonnes orthonormées sans changer le
déterminant de plus qu'un produit de longueurs.}
\begin{questions}
  \item Gram-Schmidt sur les colonnes écrit $C_j=w_j+(\text{combinaison
    des précédentes})$ avec $(w_j)$ orthogonaux : les opérations sur
    colonnes laissent $\det A=\det(w_1,\dots,w_n)$, et pour des colonnes
    orthogonales le déterminant vaut (au signe près) le produit des
    normes. Comme $\norme{w_j}\leqslant\norme{C_j}$ (le $w_j$ est une
    projection orthogonale de $C_j$) :
    $\res{|\det A|\leqslant\norme{C_1}\cdots\norme{C_n}}$.
  \item Égalité si et seulement si les colonnes sont deux à deux
    \res{\text{orthogonales}} (ou l'une nulle). Diagonale :
    $|{\det}|=|2\times3\times1|$, égalité. Pour l'exemple :
    $|\det|=3\leqslant\sqrt2\,\sqrt5\,\sqrt2=2\sqrt5\approx4{,}47$ :
    inégalité stricte, colonnes non orthogonales.
  \item Le volume d'un parallélépipède penché est au plus le produit des
    longueurs de ses arêtes — atteint quand elles sont perpendiculaires
    (chapitre 14).
\end{questions}
\end{corrige}
\end{blocLivre}

\begin{blocLivre}{cor-18-24}
\begin{corrige}{24}{Les racines des polynômes orthogonaux}
\demarche{Orthogonal à tous les petits degrés : l'argument de signes du
chapitre 15 s'applique mot pour mot.}
\begin{questions}
  \item $\R_{n-1}[X]=\operatorname{Vect}(P_0,\dots,P_{n-1})$ (degrés
    échelonnés), et $P_n$ est orthogonal à chacun : par linéarité,
    $\displaystyle\int_0^1P_n\,Q=0$ pour tout $Q\in\R_{n-1}[X]$.
  \item L'exercice 22 du chapitre 15 (avec $n-1$ au lieu de $n$) donne
    au moins $n$ changements de signe de $P_n$ dans $\intOO{0}{1}$,
    donc au moins $n$ racines distinctes ; le degré $n$ en interdit
    davantage : \res{\text{exactement }n\text{ racines réelles simples,
    toutes dans }\intOO{0}{1}}.
  \item $P_2$ : racines $\dfrac{3\pm\sqrt3}{6}\approx0{,}21$ et
    $0{,}79$ ; $P_3$ : trois racines réelles
    $\approx0{,}11$, $0{,}50$, $0{,}89$ — toutes dans l'intervalle,
    comme promis. Ces racines servent de nœuds optimaux d'intégration
    numérique : la boucle est bouclée avec les chapitres 10, 15 et 16.
\end{questions}
\end{corrige}
\end{blocLivre}

% !TEX root = ../../main.tex
%% Corrigé du défi du chapitre 18 — vérifié par verifications/defis_16_20.py.

\begin{blocLivre}{cor-18-25a}
\begin{corrige}{25}{Séries de Fourier et inégalité de Bessel}
\demarche{Linéariser pour l'orthonormalité ; projeter sur une base
orthonormée ; Bessel et Pythagore traduisent la géométrie en une inégalité
sur $\sum 1/k^{2}$.}
\partieExo{A}{Une famille orthonormée}
\begin{questions}
  \item Bilinéaire, symétrique, positif ; défini car $\int_{-\pi}^{\pi}f^{2}=0$
    impose $f=0$ pour $f$ continue (chapitre 15) : \res{\text{produit
    scalaire}}.
  \item $\cos kx\cos lx=\frac12\bigl(\cos(k-l)x+\cos(k+l)x\bigr)$ et
    $\int_{-\pi}^{\pi}\cos mx\,\mathrm{d}x=0$ pour $m\neq0$ ($2\pi$ pour
    $m=0$) : $\langle c_k,c_l\rangle=\delta_{kl}$ ($k,l\geqslant1$). De même
    $\sin kx\sin lx=\frac12\bigl(\cos(k-l)x-\cos(k+l)x\bigr)$ donne
    $\langle s_k,s_l\rangle=\delta_{kl}$ ; $\cos kx\sin lx$ est impaire,
    d'intégrale nulle ; enfin $\langle c_0,c_0\rangle=\frac1\pi\cdot\frac12\cdot2\pi=1$
    et $\langle c_0,c_k\rangle=\langle c_0,s_k\rangle=0$. La famille est
    \res{\text{orthonormée}}.
  \item En base orthonormée, $p_{T_n}(f)=\sum_e\langle f,e\rangle e$. Le
    terme constant vaut $\langle f,c_0\rangle c_0=\frac12\cdot\frac1\pi\int f=\frac{a_0}{2}$,
    et les autres $\langle f,c_k\rangle c_k=a_k\cos kx$,
    $\langle f,s_k\rangle s_k=b_k\sin kx$ : c'est $\res{S_n(f)}$.
\end{questions}
\partieExo{B}{Le cas de $f(x)=x$}
\begin{questions}[resume]
  \item $x\cos kx$ est impaire : $a_k=0$. Par parité puis intégration par
    parties,
    $b_k=\frac2\pi\int_0^{\pi}x\sin kx\,\mathrm{d}x
    =\frac2\pi\Bigl(\Bigl[-\frac{x\cos kx}{k}\Bigr]_0^{\pi}+\int_0^{\pi}\frac{\cos kx}{k}\,\mathrm{d}x\Bigr)
    =\frac2\pi\cdot\frac{-\pi(-1)^{k}}{k}=\res{\frac{2(-1)^{k+1}}{k}}$.
  \item $\norme f^{2}=\frac1\pi\int_{-\pi}^{\pi}x^{2}\,\mathrm{d}x=\frac{2\pi^{2}}{3}$.
    Bessel : $\sum_{k=1}^{n}b_k^{2}=\sum_{k=1}^{n}\frac4{k^{2}}\leqslant\frac{2\pi^{2}}3$,
    soit $\res{\sum_{k=1}^{n}\frac1{k^{2}}\leqslant\frac{\pi^{2}}6}$. Les
    sommes partielles, croissantes, sont majorées : la série converge et
    $\zeta(2)\leqslant\frac{\pi^{2}}6$.
  \item Pythagore : $\norme f^{2}=\norme{S_n(f)}^{2}+d(f,T_n)^{2}$, d'où
    $d(f,T_n)^{2}=\frac{2\pi^{2}}3-4\sum_{k\leqslant n}\frac1{k^{2}}$. Si
    $d(f,T_n)\to0$, alors $4\zeta(2)=\frac{2\pi^{2}}3$ :
    $\res{\zeta(2)=\frac{\pi^{2}}{6}}$.
\end{questions}
\end{corrige}
\end{blocLivre}

\begin{blocLivre}{cor-18-25b}
\begin{corrigeSuite}
\partieExo{C}{Convergence ponctuelle et phénomène de Gibbs}
\begin{questions}[start=7]
  \item $\sin\frac{k\pi}2$ est nul pour $k$ pair et vaut $(-1)^{j}$ pour
    $k=2j+1$ ; comme $(-1)^{k+1}=1$ pour $k$ impair :
    $S_n(f)\bigl(\frac\pi2\bigr)=2\sum_j\frac{(-1)^{j}}{2j+1}$. Si cette
    quantité tend vers $\frac\pi2$ :
    $\res{1-\tfrac13+\tfrac15-\tfrac17+\cdots=\tfrac\pi4}$ — la série de
    Leibniz, retrouvée par Fourier.
  \item $\sin k\pi=0$ : $S_n(f)(\pi)=0$ pour tout $n$, moyenne des « valeurs
    au bord » $\pi$ et $-\pi$ du prolongement $2\pi$-périodique de $f$, qui
    saute de $2\pi$. Près de $\pi$, $S_n(f)$ dépasse $f$ d'environ
    $0{,}56$, soit $9\,\%$ du saut : c'est le \emph{phénomène de Gibbs}. Il
    n'y a pas de contradiction : $d(f,T_n)$ mesure un écart \emph{moyen}
    (quadratique) ; la zone où $S_n(f)$ s'écarte de $f$ a une largeur de
    l'ordre de $\frac1n$, donc un poids qui tend vers $0$, même si l'écart
    maximal, lui, ne tend pas vers $0$. Convergence en moyenne
    quadratique ne signifie pas convergence uniforme.
\end{questions}
\end{corrigeSuite}
\end{blocLivre}

